符号距离场局部几何前馈为在不改写底层 LCP 求解主链的前提下增强互补式刚体接触建模提供了一条低侵入路径。基于速度级时间步进法（time-stepping schemes）与线性互补问题（Linear Complementarity Problem, LCP）的非平滑接触动力学，仍然是机器人仿真、计算机图形学及工程多体动力学中处理刚体碰撞、间歇接触与状态切换的核心架构之一\cite{flores2022contact,wang2021nonsmooth,stewart1996implicit,anitescu1997formulating}。在这类方法中，互补约束能够较自然地处理非穿透条件、冲击响应与接触状态切换，因此在刚体间歇接触问题上具有明确优势。与此同时，接触建模的实际表现并不只取决于互补约束的代数形式，也显著依赖于局部几何信息进入模型的方式。

传统的接触检测通常依赖于多边形网格（polygon mesh）及基于边界距离的碰撞查询，并以 GJK 及其后续扩展为代表\cite{gilbert1988fast,wang2021collision}。近年来，随着稀疏体素存储、GPU 友好数据结构与 SDF 碰撞处理的发展，利用有向距离场（Signed Distance Field, SDF）作为刚体表面的隐式代理（implicit surrogate）成为了新的热点\cite{museth2021nanovdb,macklin2020sdfcollision,liu2024sdfsdf,lai2024unifying,lopez2024convexsdf}。然而，在互补式刚体接触框架中，若局部几何仅通过距离值 $\Phi$ 与一阶梯度 $\nabla \Phi$ 进入模型，则在接触建立、持续滑移以及局部状态切换过程中，接触法向、间隙演化与约束响应往往会对局部几何误差表现出较强敏感性。对于凸轮尖端、滚子过渡区和齿轮边界等局部几何变化明显的场景，这种敏感性会进一步放大为速度毛刺、峰值误差和贴合行为退化。

在这一背景下，本文的研究焦点不再放在底层互补求解器的重新设计上，而是放在 SDF 几何代理进入既有求解链时的输入层建模问题上。对于许多工程算例，决定接触响应质量的关键因素往往已转化为局部几何信息能否以稳定且足够准确的方式进入接触模型，尤其体现在 gap 预测、局部法向流形描述以及接触建立阶段的几何噪声控制上。由此出发，相比重新构造新的刚体接触求解主链，在既有 LCP 框架内以低侵入方式增强 geometry-to-solver interface，更符合本文所处理的问题结构，也更贴近实际工程实现路径。

需要强调的是，本文并不试图让互补式方法去覆盖罚函数方法更擅长的软接触或复杂分布式接触问题。本文所关注的是互补式刚体接触框架适用范围内一个更基础的建模问题：\textbf{在不改变现有硬接触互补求解主框架、也不转向罚函数软接触建模的前提下，局部接触几何应如何表示，才能在接触建立、持续接触与状态切换过程中，以可控附加代价保持计算效率并提升接触响应精度。}换言之，本文希望回答的不是某一种增强几何代理是否可以统一优于传统方法，而是局部几何信息如何以与接触状态相匹配的方式进入互补式刚体接触模型，并在明确适用边界内转化为兼顾计算效率与响应精度的数值收益。

针对上述问题，本文提出一种\textbf{基于局部几何感知的互补式刚体接触建模思路}。在统一的 Moreau--Jean 时步框架内，本文利用 SDF 所提供的局部几何信息构造接触区附近的局部表示，并依据接触状态进行 regime 配置。对于局部接触演化连续、滑移过程较平稳的情形，采用带有 persistent manifold 的 \texttt{SlidingPatch} 表示，并在该局部流形上注入基于 Hessian 的曲率前馈；对于接触建立、局部接触域紧致或状态切换敏感的情形，则采用更紧致、更保守的 \texttt{CompactDiscrete} 表示，以保持求解稳健性。当前实现中，这一选择通过算例或应用级的预设配置完成，而不是在运行时由统一分类器自动切换。由此，本文关注的重点不是替换底层求解器，也不是提出一个自动 regime 判别算法，而是在既有互补式刚体接触框架内分析：当局部表示与接触状态相匹配时，局部几何信息能否以有限附加开销在保持计算效率的同时改善数值精度。

本工作的主要贡献包括：
\begin{itemize}
    \item \textbf{提出基于局部几何感知的互补式刚体接触建模思路}：在现有且较为成熟的互补式刚体接触求解框架内，不替换全局求解器，而是在接触建模层面引入局部 SDF 几何信息，用于增强接触区附近的几何表达能力，从而把改进重点放在 geometry-to-solver interface 上，并为在适用场景下同时兼顾计算效率与数值精度提供建模基础。
    \item \textbf{构建面向接触状态的局部表示配置方式}：对于连续滑移主导的接触，采用连续型 \texttt{SlidingPatch} 表示并结合局部曲率前馈；对于接触建立、局部紧致接触及状态切换敏感的情形，采用紧致离散的 \texttt{CompactDiscrete} 表示，使局部几何信息的使用建立在与接触 regime 相匹配的表示基础之上。
    \item \textbf{基于现有算例验证局部几何感知建模的适用性}：依托对心碰撞、偏心圆盘--滚子解析基准、Onset Stress Benchmark、工业凸轮和 Twin Gear 等算例，系统评估了该建模思路在不同接触场景中的表现。结果表明，在与接触状态相匹配的 regime 下，该方法能够在连续滑移与接触建立问题中改善局部几何描述与整体数值响应，并体现出在保持计算效率的同时改善数值精度的特征；而在局部奇异边界或非光滑啮合场景中，其有效性更依赖于局部表示方式与接触状态之间的匹配。
\end{itemize}

需要进一步说明的是，本文实验并不以证明某一种单一 regime 对全部 benchmark 普遍最优为目标。相反，本文将各算例视为不同接触状态下的代表性场景，并在与其局部几何特征相匹配的预设 regime 下考察该建模思路是否能够提供稳定、可解释且兼顾计算效率与数值精度的收益。因此，本文最终建立的是一种面向效率与精度协同兼顾的建模原则，而不是一套运行时自动选择最优局部表示的统一算法。
